% ch3.tex — Hermite Polynomials and Applications
% Hand-typeset精修版，对应原书 Chapter 3 (book pages 101–132)

\chapter{Hermite Polynomials and Applications}

\section{Introduction}
\label{sec:3.1}

In this chapter and the next, we will discuss the theory of orthogonal
polynomials and their applications to some problems in Physics. We will
essentially discuss the case of Hermite and Laguerre polynomials (HP,
LP) and touch on the other families.

In the previous chapters, we have ``evocated'' the Hermite type
polynomials many times and we have shown that they can be defined
through the operational identity (see eq.~(2.2.19))
\begin{equation}
H_n(x,y)=e^{y\partial_x^2}(x^n).
\label{eq:3.1.1}
\end{equation}
Even though we have already left as exercise the explicit derivation of
$H_n(x,y)$, we give it here for the sake of completeness, in view of
their importance for the present book.

By expanding the exponential operator, we find
\begin{equation}
H_n(x,y)=\sum_{s=0}^{\infty}\frac{y^s}{s!}\,\partial_x^{2s}(x^n).
\label{eq:3.1.2}
\end{equation}

\noindent{\itshape Remark.} The polynomials $H_n(x,y)$ are not the Hermite
polynomials in their canonical form. They are sometimes called heat,
Kamp\'e de F\'eriet, Gould--Hopper \ldots{} polynomials. In the
following we will refer to them as Hermite or heat polynomials.

By using the identity
\[
\partial_x^m(x^n)=
\begin{cases}
\dfrac{n!}{(n-m)!}\,x^{n-m}, & n>m,\\[6pt]
0, & n<m,
\end{cases}
\label{eq:3.1.3}
\]
we obtain
\begin{equation}
H_n(x,y)=n!\sum_{s=0}^{\lfloor n/2\rfloor}
\frac{x^{\,n-2s}\,y^s}{(n-2s)!\,s!},
\qquad H_n(x,0)=x^n.
\label{eq:3.1.4}
\end{equation}

The above family of polynomials are not the HP in their canonical
form but belong to the Hermite-like polynomial families. As already
mentioned, they are known as the \emph{heat polynomials} since they
provide a natural solution to the heat equation and, as we will show in
the following, they are not orthogonal for all values of the $y$
coordinate.

In Fig.~\ref{fig:3.1} we have reported the first four heat
polynomials~\eqref{eq:3.1.4} vs.\ the variable $x$ for two different
values of $y$ (playing the role of a parameter).

\begin{figure}[H]
\centering
\includegraphics[max width=0.95\linewidth,keepaspectratio]{fig3_1.png}\caption{Heat polynomials $H_n(x,y)$ vs.\ the variable $x$ for
(a)~$y=-1$ and (b)~$y=+1$.}
\label{fig:3.1}
\end{figure}

We can note a first substantial difference: the polynomials with
negative $y$ can take positive and negative values and have a number
of zeros equal to the order of the polynomial. In the other case, even
order polynomials have no real zeros and the odd order are
anti-symmetric with only one zero at the origin.

\medskip
The operational definition~\eqref{eq:3.1.1} greatly simplifies the
study of the properties of HP and of their generalizations as well.
From this definition, it is evident that
\[
H_{n+1}(x,y)=e^{y\partial_x^2}(x^{n+1})
=e^{y\partial_x^2}\,x\,e^{-y\partial_x^2}\,e^{y\partial_x^2}(x^n)
=e^{y\partial_x^2}\,x\,e^{-y\partial_x^2}\,H_n(x,y).
\]
By taking into account eq.~\eqref{eq:2.4.16},
\[
e^{y\partial_x^2}\,x\,e^{-y\partial_x^2}=x+2y\partial_x,
\]
we obtain the recurrence relation
\begin{equation}
H_{n+1}(x,y)=x\,H_n(x,y)+2y\,\partial_x H_n(x,y).
\label{eq:3.1.7}
\end{equation}
The use of an analogous procedure yields the recurrence relation
\begin{equation}
\partial_x H_n(x,y)=n\,H_{n-1}(x,y).
\label{eq:3.1.8}
\end{equation}
These two recurrences can be combined to prove that the heat
polynomials satisfy the ODE
\begin{equation}
2y\,z_{xx}+x\,z_x=n\,z,\qquad z=H_n(x,y).
\label{eq:3.1.9}
\end{equation}
In the above equation, the derivative taken with respect to $x$ and
$y$ is viewed as a parameter.

It is also easily checked that, by considering $y$ as a variable,
the following relation holds:
\begin{equation}
\partial_y H_n(x,y)=n(n-1)\,H_{n-2}(x,y).
\label{eq:3.1.10}
\end{equation}

To frame the previous results in a more general context, we introduce
the notion of \emph{quasi-monomials}: those families of polynomials
$p_n(x)$ for which one can define a couple of operators $\hat P$ and
$\hat M$ (referred to as derivative and multiplicative operators,
respectively) such that
\begin{equation}
\hat P\,p_n(x)=n\,p_{n-1}(x),\qquad
\hat M\,p_n(x)=p_{n+1}(x).
\label{eq:3.1.11}
\end{equation}
Furthermore, if we add the condition
\begin{equation}
p_0(x)=1,
\label{eq:3.1.12}
\end{equation}
we can construct the quasi-monomial family using the multiplicative
operator only, according to the identity
\begin{equation}
\hat M^n\,1=p_n(x).
\label{eq:3.1.13}
\end{equation}
which is obtained by iteration from the second of
eq.~\eqref{eq:3.1.11}. Furthermore, if the operators $\hat M$ and
$\hat P$ have a differential realization, the polynomials $p_n(x)$ are
easily shown to satisfy the differential equation
\begin{equation}
\hat M\hat P\,p_n(x)=n\,p_n(x).
\label{eq:3.1.14}
\end{equation}

In the following, we will discuss a general procedure to construct
quasi-monomials starting from the ordinary monomial $x^n$, whose
relevant multiplicative and derivative operators are $x$ and the
ordinary derivative, respectively. It is evident that, according to
the above definition, the $H_n(x,y)$ are quasi-monomials and that
the relevant multiplicative and derivative operators can be identified as
\begin{equation}
\hat M=x+2y\partial_x,\qquad \hat P=\partial_x
\label{eq:3.1.15}
\end{equation}
and
\begin{equation}
\hat M\hat P=x\partial_x+2y\partial_x^2.
\label{eq:3.1.16}
\end{equation}
Therefore, HP are eigenfunctions of the operator~\eqref{eq:3.1.16}
with (real) eigenvalues~$n$.

\noindent{\itshape Remark.} In this case $y$ should be considered a
parameter and not a variable. The notation of total derivative
should be used but we use the partial derivative for continuity of
notation.

The above introductory remarks provide the back-bone of the
discussion to be developed in the forthcoming sections.

% ------------------------------------------------------------------
\section{Hermite Polynomials Generating Function}
\label{sec:3.2}

\subsection{Introducing the Generating Function}
\label{sec:3.2.1}

Being the polynomials $H_n(x,y)$ solutions of the heat equation,
they can be introduced through an operational definition involving an
evolution operator with a second order derivative as argument. The
following two important identities are a direct consequence of such a
definition:
\begin{equation}
e^{z\partial_x^2}H_n(x,y)=H_n(x,y+z),\qquad
e^{-y\partial_x^2}H_n(x,y)=x^n.
\label{eq:3.2.1}
\end{equation}
which, albeit straightforward, are very important and should be
carefully considered.

A further important property of Hermite (heat) polynomials is the
generating function, which will be widely exploited in the following.
By generating function associated with a given polynomial family
$p_n(x)$, we mean the function
\begin{equation}
G(x,t)=\sum_{n=0}^{\infty}\frac{t^n}{n!}\,p_n(x)
\label{eq:3.2.2}
\end{equation}
which is essentially a function of the variables $t$ and $x$. The
relevant expansion in a Taylor series in the variable $t$, called the
expansion variable of the generating function, yields the
coefficients $\partial_t^n G(x,t)|_{t=0}$ which are just identified
with $p_n(x)$ polynomials. It is evident that the ordinary
monomials are generated by
\begin{equation}
\sum_{n=0}^{\infty}\frac{t^n}{n!}\,x^n=e^{xt}.
\label{eq:3.2.3}
\end{equation}

According to eq.~\eqref{eq:3.1.1}, we can define the generating
function of the HP as
\[
\sum_{n=0}^{\infty}\frac{t^n}{n!}\,H_n(x,y)
=e^{y\partial_x^2}\!\left(\sum_{n=0}^{\infty}\frac{t^n}{n!}\,x^n\right)
=e^{y\partial_x^2}\,e^{xt}.
\]
By using $e^{y\partial_x^2}\,e^{xt}=e^{xt+yt^2}$ (since
$\partial_x^m e^{xt}=t^m e^{xt}$), we get
\begin{equation}
\sum_{n=0}^{\infty}\frac{t^n}{n!}\,H_n(x,y)=e^{xt+yt^2}.
\label{eq:3.2.5}
\end{equation}

\noindent{\itshape Remark.} It is to be noted that
$e^{2yt\partial_x}\neq 1$ in general, hence the intermediate steps
require the operator disentanglement
$e^{y\partial_x^2}\,e^{xt}=e^{yt^2}\,e^{xt}\,e^{2yt\partial_x}$ and
subsequent action on the vacuum $1$.

\medskip
As will be shown in the following, this is an extremely important
relation in the theory of Hermite polynomials.

We have derived the generating function of HP as a by-product of an
operational formalism and after having stated the main properties of
this polynomial family. In more classical textbooks on special
functions, the explicit form of the HP is inferred directly from the
generating function. To understand such a procedure, we remind the
Cauchy formula for the series product: namely, given the two infinite
series
\[
S_1=\sum_{s=0}^{\infty}b_s\,x^s,\qquad
S_2=\sum_{r=0}^{\infty}a_r\,x^r,
\]
we obtain
\[
S_1S_2=\sum_{s=0}^{\infty}\sum_{r=0}^{\infty}b_s\,a_r\,x^{r+s}
=\sum_{n=0}^{\infty}c_n\,x^n,\qquad
c_n=\sum_{r=0}^{n}a_r\,b_{n-r}.
\]

Therefore, we can directly derive the HP from the generating
function~\eqref{eq:3.2.5}:
\[
e^{xt+yt^2}=e^{xt}\,e^{yt^2}
=\sum_{r=0}^{\infty}\frac{x^r t^r}{r!}
 \sum_{s=0}^{\infty}\frac{y^s t^{2s}}{s!}
=\sum_{n=0}^{\infty}\frac{t^n}{n!}\,H_n(x,y),
\]
and, by setting $r+2s=n$, according to the
definition~\eqref{eq:3.1.4}, we get
\[
\sum_{n=0}^{\infty}\frac{t^n}{n!}\,H_n(x,y)
=e^{xt+yt^2}.
\]

\subsection{Generating Function Applications}
\label{sec:3.2.2}

Integrals of the type
\begin{equation}
I_n(a,b)=\int_{-\infty}^{\infty}(x+b)^n\,e^{-ax^2}\,dx
\label{eq:3.2.11}
\end{equation}
occur very often in the solution of problems in Physics or
Engineering. The use of the generating function method is extremely
efficient to evaluate these integrals. If we note that
\[
\sum_{n=0}^{\infty}\frac{t^n}{n!}\,I_n(a,b)
=\int_{-\infty}^{\infty}e^{tb}\,e^{tx}\,e^{-ax^2}\,dx
=e^{tb}\,\sqrt{\frac{\pi}{a}}\;e^{\frac{t^2}{4a}},
\label{eq:3.2.12}
\]
according to eq.~\eqref{eq:3.2.5} and by applying the principle of
identity of polynomials to eq.~\eqref{eq:3.2.12}, namely by equating
the like $t$-powers, we obtain the following result
\begin{equation}
I_n(a,b)=\sqrt{\frac{\pi}{a}}\;\frac{1}{a^{n/2}}\,
H_n\!\left(\frac{b}{\sqrt{a}},\frac{1}{4a}\right).
\label{eq:3.2.13}
\end{equation}

The further two properties
\begin{equation}
\partial_x^n e^{ax^2}=H_n(2ax,a)\,e^{ax^2},
\label{eq:3.2.14}
\end{equation}
\begin{equation}
a^n H_n(x,y)=H_n(ax,a^2y),
\label{eq:3.2.15}
\end{equation}
can be proved by means of the generating function method.
Eq.~\eqref{eq:3.2.14} is indeed easily proved as a consequence of
the properties of the shift operators:
\[
\sum_{n=0}^{\infty}\frac{t^n}{n!}\,\partial_x^n e^{ax^2}
=e^{t\partial_x}e^{ax^2}=e^{a(x+t)^2}
=e^{ax^2}\,e^{2atx+at^2}
=e^{ax^2}\sum_{n=0}^{\infty}\frac{t^n}{n!}\,H_n(2ax,a),
\]
while eq.~\eqref{eq:3.2.15} is left as exercise.

\noindent{\itshape Remark.} Note that an important consequence of
eq.~\eqref{eq:3.2.15} is the identity
$(-1)^n H_n(x,y)=H_n(-x,y)$.

\medskip
As a further example of the application of the generating function
method, we consider the proof of the following Theorem.

\begin{theorem}[Addition]
$\forall\,x,y,z\in\mathbb{R},\;\forall n\in\mathbb{N}$
\begin{equation}
H_n(x+z,y)=\sum_{s=0}^{n}\binom{n}{s}\,
H_{n-s}\!\left(x,\frac{y}{2}\right)
H_s\!\left(z,\frac{y}{2}\right).
\label{eq:3.2.17}
\end{equation}
\end{theorem}

\noindent{\itshape Hint:} it is just a consequence of the fact that
\[
\sum_{n=0}^{\infty}\frac{t^n}{n!}\,H_n(x+z,y)
=e^{(x+z)t+yt^2}
=e^{xt+yt^2/2}\;e^{zt+yt^2/2}.
\tag{3.2.18}
\]

The following problem is extremely important, since it deals with a
non-trivial ordering problem and yields an idea of how the HP appear,
in a quite natural way, in manipulating expressions containing
multiplicative derivative operators. The problem consists in finding
the analogue of the Newton binomial form for the operator
\begin{equation}
\hat O_n=(\alpha\partial_x+\beta x)^n.
\label{eq:3.2.19}
\end{equation}

According to the previous discussion, we have
\[
\sum_{n=0}^{\infty}\frac{t^n}{n!}\,\hat O_n
=e^{t(\alpha\partial_x+\beta x)}
=e^{\beta t x+\frac{\alpha\beta}{2}t^2}\;e^{\alpha t\partial_x},
\]
and, taking into account eq.~\eqref{eq:3.2.5}, we can write
\[
e^{\beta t x+\frac{\alpha\beta}{2}t^2}
=\sum_{r=0}^{\infty}\frac{t^r}{r!}\,
H_r\!\left(\beta x,\frac{\alpha\beta}{2}\right)
=\sum_{r=0}^{\infty}\frac{t^r}{r!}\,H_r\!\left(\beta x,\frac{\alpha\beta}{2}\right)
 \sum_{s=0}^{\infty}\frac{(\alpha t)^s}{s!}\,\partial_x^s,
\]
which, after applying the Cauchy product formula and comparing with
$\sum t^n\hat O_n/n!$, yields
\begin{equation}
\hat O_n=\sum_{s=0}^{n}\frac{1}{s!}\,
\alpha^s\,H_{n-s}\!\left(\beta x,\frac{\alpha\beta}{2}\right)\,\partial_x^s,
\label{eq:3.2.22}
\end{equation}
which is a generalization of the so-called \emph{Burchnall}
operational rule and can be used to prove the identity
\begin{equation}
\hat O_n\,e^{-x^2}
=\sum_{s=0}^{n}\frac{(-\alpha)^s}{s!}\,
H_{n-s}\!\left(\beta x,\frac{\alpha\beta}{2}\right)
H_s(2x,-1)\,e^{-x^2}
=H_n\!\left((\beta-2\alpha)x,\frac{\alpha(\beta-2\alpha)}{2}\right)e^{-x^2}.
\label{eq:3.2.23}
\end{equation}

We will discuss further examples of applications of the generating
function method in the second part of this book.

% ------------------------------------------------------------------
\section{Hermite Polynomials as an Orthogonal Basis}
\label{sec:3.3}

We will more carefully define the concept of orthogonal polynomials
later. Here we will address the following problem: under which
condition can a given function of $x$ be expanded in terms of
Hermite polynomials?

To this aim, we treat the second variable $y$ as a parameter and
limit ourselves to the case $y<0$. We therefore write
\begin{equation}
F(x)=\sum_{n=0}^{\infty}a_n\,H_n(x,-|y|).
\label{eq:3.3.1}
\end{equation}
Our task is now to evaluate the coefficients $a_n$ of the above
expansion. According to the definition~\eqref{eq:3.1.1}, we can
recast~\eqref{eq:3.3.1} in the form
\[
F(x)=e^{-|y|\partial_x^2}\!\left(\sum_{n=0}^{\infty}a_n\,x^n\right),
\]
which can be inverted to get
\[
\Phi(x)\equiv e^{|y|\partial_x^2}\!\bigl(F(x)\bigr)
=\sum_{n=0}^{\infty}a_n\,x^n.
\]
According to the Gauss--Weierstrass transform (see
eq.~\eqref{eq:2.2.17}), the following integral representation for
the function $\Phi(x)$ holds:
\begin{equation}
\Phi(x)=\frac{1}{\sqrt{4\pi|y|}}\int_{-\infty}^{\infty}
e^{-\frac{(x-\sigma)^2}{4|y|}}\,F(\sigma)\,d\sigma.
\label{eq:3.3.4}
\end{equation}
By using the generating function of HP, the last identity can be
written as
\[
\Phi(x)=\frac{1}{\sqrt{4\pi|y|}}
 \sum_{n=0}^{\infty}\frac{1}{n!}
 \int_{-\infty}^{\infty}
 H_n\!\left(\frac{\sigma}{2\sqrt{|y|}},-\frac{1}{4}\right)
 e^{-\frac{(x-\sigma)^2}{4|y|}}\,F(\sigma)\,d\sigma,
\]
which, compared to $\Phi(x)=\sum a_n x^n$, gives for the
coefficients $a_n$ the expression
\begin{equation}
a_n=\frac{1}{2^n n!\sqrt{\pi|y|}}
\int_{-\infty}^{\infty}
H_n\!\left(\frac{\sigma}{2\sqrt{|y|}},-\frac{1}{4}\right)
e^{-\frac{\sigma^2}{4|y|}}\,F(\sigma)\,d\sigma.
\label{eq:3.3.6}
\end{equation}

We remark that, as a consequence of eq.~\eqref{eq:3.3.1}, the
functions
\begin{equation}
\phi_n(x)=\frac{1}{\sqrt{2^n n!\,\pi|y|}}\,
H_n\!\left(\frac{x}{2\sqrt{|y|}},-\frac{1}{4}\right)
e^{-\frac{x^2}{4|y|}}
\label{eq:3.3.7}
\end{equation}
are bi-orthogonal to the Hermite polynomials
$H_n(x,-|y|)$. In more conventional terms, we can state that the
polynomials $H_n\!\left(\frac{x}{2\sqrt{|y|}},-\frac{1}{4}\right)$ are
orthogonal to the $H_n(x,-|y|)$ with weight function
$e^{-x^2/(4|y|)}$. The canonical form of the HP uses $|y|=1/2$ and
is denoted by $\mathrm{He}_n(x)$, i.e.
\begin{equation}
\mathrm{He}_n(x)=H_n\!\left(x,-\frac{1}{2}\right).
\label{eq:3.3.8}
\end{equation}
In the applications, the form
\begin{equation}
H_n(x)=H_n(2x,-1)=2^n H_n\!\left(x,-\frac{1}{4}\right)
\label{eq:3.3.9}
\end{equation}
is also used and we will adopt both notations. These two forms are
linked by the identity
\begin{equation}
H_n(x)=2^{n/2}\,\mathrm{He}_n\!\left(\frac{x}{\sqrt{2}}\right).
\label{eq:3.3.10}
\end{equation}
Furthermore, the link of the $H_n(x,y)$ to the Hermite canonical
form is provided by
\begin{equation}
H_n(x,y)=(-i)^n\,y^{n/2}\,
H_n\!\left(\frac{x}{2\sqrt{-y}}\right)
=2^n\,(-iy)^{n/2}\,
\mathrm{He}_n\!\left(\frac{x}{2\sqrt{-iy}}\right).
\label{eq:3.3.11}
\end{equation}

As it can be easily checked, the above relations are a
straightforward consequence of the previously outlined properties.

\noindent{\itshape Remark.} Two sets of functions $\gamma_n(x)$,
$\lambda_n(x)$ are said \emph{bi-orthogonal} if
$\int \gamma_n(x)\,\lambda_m(x)\,dx=\delta_{nm}$. For example
$\sin(nx)$, $\cos(mx)$ are bi-orthogonal in the interval
$[0,2\pi]$.

\medskip
Let us now derive the differential equation satisfied by the
function~\eqref{eq:3.3.7}. It will be determined by means of a
simple but effective procedure. From eq.~\eqref{eq:3.3.7}, one has
\[
H_n\!\left(\frac{x}{2\sqrt{|y|}},-\frac{1}{4}\right)
=A_n\,e^{\frac{x^2}{4|y|}}\,\phi_n(x),\qquad
A_n=\sqrt{2^n n!\,\pi|y|}.
\]

In operational terms, the ODE specifying the differential equation of
the polynomials $H_n(x,y)$ can be written as (see
eqs.~\eqref{eq:3.1.14} and~\eqref{eq:3.1.16})
\[
\hat h_{\mathrm{er}}\,H_n(ax,y)=0,\qquad
\hat h_{\mathrm{er}}=\partial_x^2+a^2x\partial_x-n a^2.
\]
The operator $\hat h_{\mathrm{er}}$ applied to eq.~\eqref{eq:3.3.7}
yields
\[
\left(-2|y|\partial_x^2+x\partial_x-n\right)e^{\frac{x^2}{4|y|}}\phi_n(x)=0,
\]
and thus, for $\phi_n(x)$, we obtain the second order differential
equation with non-constant coefficients
\begin{equation}
\left(-2|y|\partial_x^2-x\partial_x-1\right)\phi_n(x)=n\,\phi_n(x).
\label{eq:3.3.15}
\end{equation}
This result means that the function $\phi_n$ is the eigenfunction
with eigenvalue $n$ of the operator
\begin{equation}
\hat\eta_{\mathrm{her}}=-2|y|\partial_x^2-x\partial_x-1,
\label{eq:3.3.16}
\end{equation}
which, for $y<0$, is the Hermitian conjugate of the
operator in eq.~\eqref{eq:3.1.16}.

\noindent{\itshape Remark.} We remind that the Hermitian conjugate of
the product of two operators is
$(\hat A\hat B)^\dagger=\hat B^\dagger\hat A^\dagger$ and being
$\hat x^\dagger=\hat x$ and $\partial_x^\dagger=-\partial_x$, we find
$(x\partial_x)^\dagger=-\partial_x x\ldots$

\medskip
Even though without proof, we underscore the following statement.

\begin{proposition}
The eigenfunctions of two operators Hermitian conjugate to each
other form a bi-orthogonal set.
\end{proposition}

Let us now consider the case of the canonical Hermite polynomials
with $y=-1/2$ and define the function
\begin{equation}
\Phi_n(x)=\alpha_n\,\mathrm{He}_n(x)\,e^{-x^2/4}.
\label{eq:3.3.17}
\end{equation}
which, evidently, provides an orthogonal set of functions whose
normalization constant $\alpha_n$ will be specified later. The
differential equation satisfied by the orthogonal
functions~\eqref{eq:3.3.17} can be obtained by proceeding as
illustrated before, for the function~\eqref{eq:3.3.7}, namely, if we
get
\[
\phi_n(x)=\Phi_n(x)\,e^{x^2/4}.
\]
According to eq.~\eqref{eq:3.3.15}, the ODE satisfied by the
function $\Phi_n(x)$ is
\begin{equation}
\left(\frac{x^2}{4}-\partial_x^2-\frac{1}{2}\right)\Phi_n(x)
=n\,\Phi_n(x),
\label{eq:3.3.19}
\end{equation}
which is recognized as the quantum harmonic oscillator equation. The
harmonic oscillator eigenfunctions are usually written in the base
$H_n(x)$ and read
\begin{equation}
\Phi_n(x)=\frac{1}{\sqrt{2^n n!\sqrt{\pi}}}\,
H_n(x)\,e^{-x^2/2}.
\label{eq:3.3.20}
\end{equation}

\begin{exercise}
Prove that eq.~\eqref{eq:3.3.20} provides an orthonormal set in
the sense that $\Phi_n(x)$ are orthogonal and normalized to unity.
\end{exercise}

In Fig.~\ref{fig:3.2} we have shown the behavior of the first six
harmonic oscillator eigenfunctions.

\begin{figure}[H]
\centering
\includegraphics[max width=0.95\linewidth,keepaspectratio]{fig3_2.png}\caption{Harmonic oscillator eigenfunctions $\Phi_n(x)$.}
\label{fig:3.2}
\end{figure}

The expansion of functions in terms of HP will be discussed in the
part devoted to exercises and complements. Here we propose a very
simple, but important, example of this kind of expansion.

\begin{exercise}
Show that
\begin{equation}
x^n=\sum_{s=0}^{\lfloor n/2\rfloor}
\frac{n!}{s!\,(n-2s)!}\,(-1)^s\,y^s\,H_{n-2s}(x,y).
\tag{3.3.21}
\end{equation}
\noindent{\itshape Hint:} use the generating function method and note
that $e^{xt}=e^{xt+yt^2-yt^2}$.
\end{exercise}

% ------------------------------------------------------------------
\section{Hermite Polynomials in Quantum Mechanics:
         Creation and Annihilation Operators}
\label{sec:3.4}

The use of the formalism of the harmonic oscillator is almost
ubiquitous in any field of Physics and it is therefore difficult to
select examples which can be considered the most representative.
Here we follow the ``tradition'' that, although not original, has the
merit of offering a fairly solid and pedagogically effective way to
approach the problem.

We will therefore consider some examples from Quantum Mechanics
and introduce our discussion starting from the Hamiltonian of a single
particle driven by a harmonic oscillator potential, namely
\begin{equation}
\hat H=\frac{\hat p^2}{2m}+\frac{1}{2}\,k\,\hat x^2.
\label{eq:3.4.1}
\end{equation}
The associated Schr\"odinger equation writes
\begin{equation}
i\hbar\,\partial_t\Psi=-\frac{\hbar^2}{2m}\,\partial_x^2\Psi
+\frac{1}{2}\,k\,x^2\Psi
\label{eq:3.4.2}
\end{equation}
and its stationary solution can be cast in the form
\begin{equation}
\Psi(x,t)=e^{-iE_n t/\hbar}\,\Phi_n(x).
\label{eq:3.4.3}
\end{equation}

\begin{figure}[H]
\centering
\includegraphics[max width=0.95\linewidth,keepaspectratio]{fig3_3.png}\caption{Energy levels, harmonic oscillator potential and the first
eight eigenfunctions vs.\ the position coordinate.}
\label{fig:3.3}
\end{figure}

which, inserted in eq.~\eqref{eq:3.4.2}, yields
\begin{equation}
E_n\,\Phi_n(x)=-\frac{\hbar^2}{2m}\,\partial_x^2\Phi_n(x)
+\frac{1}{2}\,k\,x^2\Phi_n(x),
\label{eq:3.4.4}
\end{equation}
where $\omega=\sqrt{k/m}$ is the characteristic frequency of our
problem. Eq.~\eqref{eq:3.4.4} is just the ODE satisfied by the
harmonic oscillator eigenfunctions (note that we refer to the
eigenfunctions in the base $H_n(x)$). We remark that the quantity
\begin{equation}
x_q\equiv\sqrt{\frac{\hbar}{m\omega}}
\label{eq:3.4.5}
\end{equation}
has the dimensions of a length (its physical meaning will be
discussed below) and by rewriting eq.~\eqref{eq:3.4.4} in terms of
the variable $\xi=x/x_q$, we find
\begin{equation}
E_n\,\Phi_n(\xi)=-\frac{\hbar\omega}{2}\,\partial_\xi^2\Phi_n(\xi)
+\frac{\hbar\omega}{2}\,\xi^2\Phi_n(\xi),
\label{eq:3.4.6}
\end{equation}
with the energy eigenvalues $E_n$ which can be written in terms of
the characteristic frequency as
\begin{equation}
E_n=\hbar\omega\!\left(n+\frac{1}{2}\right).
\label{eq:3.4.7}
\end{equation}

In Fig.~\ref{fig:3.3}, we have reported the quadratic potential,
the energy levels and the first eight wave functions. It is evident
that the larger the order of the excitation, the larger the spatial
extension of the distribution function. It must be stressed that the
0-order eigenfunction, i.e.\ the lowest energy level, is specified by
\[
\Phi_0(x)\propto e^{-x^2/(2x_q^2)},
\]
i.e.\ a wave packet with an r.m.s.\ proportional to $x_q$.

In the previous discussion, we have introduced the elastic constant
$k$ which is the most important physical parameter because it
specifies the oscillation frequency and, thus, the spacing between
the energy eigenvalues. We will discuss in the forthcoming section
examples involving specific applications like the case of diatomic
molecules. Here we pursue in fixing the formalism we will use in the
following, by introducing the creation--annihilation operators.

To this aim, we cast eq.~\eqref{eq:3.4.6} in the form
\begin{equation}
\left(\xi^2-\partial_\xi^2-1\right)\Phi_n(\xi)=2n\,\Phi_n(\xi),
\label{eq:3.4.9}
\end{equation}
which suggests the introduction of the operators
\begin{equation}
\hat a=\frac{1}{\sqrt{2}}\,(\xi+\partial_\xi),\qquad
\hat a^\dagger=\frac{1}{\sqrt{2}}\,(\xi-\partial_\xi).
\label{eq:3.4.10}
\end{equation}
These operators are Hermitian conjugate to each other and satisfy the
rule of commutation
\begin{equation}
[\hat a,\hat a^\dagger]=\hat 1.
\label{eq:3.4.11}
\end{equation}
and, according to the manipulations we learned in Chapter~2, we can
recast eq.~\eqref{eq:3.4.9} in the form
\begin{equation}
\hat a^\dagger\hat a\,\Phi_n(\xi)=n\,\Phi_n(\xi).
\label{eq:3.4.12}
\end{equation}

In the previous section, we have seen how to construct the
polynomials by using the multiplicative operator acting many times on
the 0-order polynomial. We will now prove that the action of the
$n$th power of the operator $\hat a^\dagger$ on what we define as
the vacuum state, represented by the lowest order harmonic
oscillator function, yields
\begin{equation}
\frac{(\hat a^\dagger)^n}{\sqrt{n!}}\,\Phi_0=\Phi_n.
\label{eq:3.4.13}
\end{equation}

To prove the correctness of the previous identity, we use the
generating function method, according to which we find
\[
\sum_{n=0}^{\infty}\frac{t^n}{n!}\,(\hat a^\dagger)^n\Phi_0(\xi)
=e^{t\hat a^\dagger}\Phi_0(\xi)
=e^{-t^2/4}\,e^{\sqrt{t/2}\,\xi}\,
 e^{-t^2/4}\,e^{\sqrt{t/2}\,\xi}\,e^{-\xi^2/2}
\]
\begin{align*}
&=e^{-t^2/2}\,e^{\sqrt{2t}\,\xi}\,e^{-\xi^2/2}
 =\sum_{m=0}^{\infty}\frac{t^m}{m!}\,
  \mathrm{He}_m\!\left(\sqrt{2}\,\xi\right)\,e^{-\xi^2/2}.
\tag{3.4.14}
\end{align*}
From the comparison between the same power coefficients, we obtain
the proof of the identity~\eqref{eq:3.4.13} which also yields
\begin{equation}
\hat a^\dagger\Phi_n=\sqrt{n+1}\,\Phi_{n+1}.
\label{eq:3.4.15}
\end{equation}
For this reason, we will refer to $\hat a^\dagger$ as the
``creation'' operator. Along the same line, the companion identity
for the ``annihilation'' operator can also be checked by the use of
the same procedure as before:
\begin{equation}
\hat a\,\Phi_n=\sqrt{n}\,\Phi_{n-1}.
\label{eq:3.4.16}
\end{equation}

We have thus shown that three operators play a fundamental role in
the quantum harmonic oscillator dynamics: the Hermitian conjugate
creation--annihilation partners (sometimes called ladder
operators) and the Hermitian operator $\hat a^\dagger\hat a$, called
the number operator.

The Dirac notation is often exploited to describe the harmonic
oscillator states. We have, indeed, the correspondence
\begin{equation}
|0\rangle\to\Phi_0(\xi),\qquad
|n\rangle\to\Phi_n(\xi)
\label{eq:3.4.17}
\end{equation}
and the orthogonality property can be expressed as
\begin{equation}
\langle m|n\rangle=\delta_{m,n}.
\label{eq:3.4.18}
\end{equation}

In eq.~\eqref{eq:3.4.17}, we have kept separate the $|0\rangle$
state which denotes the so-called vacuum state. In more rigorous
terms the correspondence between Dirac (also called $n$-states)
and the harmonic oscillator eigenfunctions should be written as
\begin{equation}
\langle \xi|n\rangle=\Phi_n(\xi).
\label{eq:3.4.19}
\end{equation}

This aspect of the problem will be more carefully discussed in the
second part of the book, with the introduction of the so-called
coherent states.

% ------------------------------------------------------------------
\section{Quantum Mechanics Applications}
\label{sec:3.5}

Unfortunately, we cannot go deeply into the application of the above
formalism to Quantum Mechanics problems and therefore we will limit
ourselves to a few examples, which will clarify the relevant power
and elegance. The first problem we are going to discuss is that of a
charged particle ruled by a harmonic potential and under the action
of a classical electric field.

The relevant Hamiltonian can be written as
\begin{equation}
\hat H=\frac{\hat p^2}{2m}+\frac{1}{2}\,k\,\hat x^2-eE\hat x,
\label{eq:3.5.1}
\end{equation}
where the electric field $E$ is constant and homogeneous, i.e.\ it is
independent of time and coordinates. The difference with the case we
have already studied is the presence of a linear potential,
accounting for the effect of the electric field.

The above problem can be easily turned into an ordinary harmonic
oscillator problem by recasting the above Hamiltonian in the form
\begin{equation}
\hat H=\frac{\hat p^2}{2m}+\frac{1}{2}\,m\omega^2
\left(\hat x-\frac{eE}{m\omega^2}\right)^2\hat 1
-\frac{(eE)^2}{2m\omega^2}\,\hat 1.
\label{eq:3.5.2}
\end{equation}
The last term in eq.~\eqref{eq:3.5.2}, multiplied by the unit
operator, is a constant with the dimension of the energy, while
\begin{equation}
x_s\equiv\frac{eE}{m\omega^2}
\label{eq:3.5.3}
\end{equation}
is a coordinate shift induced by the presence of the electric field.
Since the electric field is constant and homogeneous, the
rearrangement of the Hamiltonian in the form~\eqref{eq:3.5.2}
suggests that the introduction of the coordinate operator
\begin{equation}
\hat X=\hat x-\frac{eE}{m\omega^2}\,\hat 1
\label{eq:3.5.4}
\end{equation}
leaves unchanged the commutation relations of our problem and
\begin{equation}
[\hat p,\hat X]=-i\hbar\,\hat 1.
\label{eq:3.5.5}
\end{equation}
We can therefore conclude that
\begin{enumerate}
\item The energy eigenvalues are the same as before but with an
      additional term due to the effect of the electric field, namely
      \begin{equation}
      E_n=\hbar\omega\!\left(n+\frac{1}{2}\right)-\frac{(eE)^2}{2m\omega^2}.
      \label{eq:3.5.6}
      \end{equation}
\item The eigenfunctions remain the same but the argument should be
      replaced by
      \begin{equation}
      \xi\to\xi-\xi_0,\qquad
      \xi_0=\frac{x_s}{x_q}=\sqrt{\frac{2}{m\hbar\omega}}\;eE.
      \label{eq:3.5.7}
      \end{equation}
\end{enumerate}

The second example we discuss is that of the Landau levels,
appearing in the quantum analysis of an electron (or any charged
particle) moving in a classical magnetic field. The Hamiltonian
operator ruling such a process can be written as (cgs units)
\begin{equation}
\hat H=\frac{1}{2m}\left(\hat{\vec p}-\frac{q}{c}\,\vec A\right)^2.
\label{eq:3.5.8}
\end{equation}
where $\vec A$ is the vector potential. By assuming that the
magnetic field is directed along the $z$ direction and that the
electron motion occurs in the $x,y$ plane,
\begin{equation}
\vec B=\nabla\times\vec A,
\label{eq:3.5.9}
\end{equation}
we can write $\vec A$ as
\begin{equation}
\vec A\equiv(0,Bx,0).
\label{eq:3.5.10}
\end{equation}
The Hamiltonian can therefore be written as
\begin{equation}
\hat H=\frac{1}{2m}\left[\hat p_x^2
+\left(\hat p_y-\frac{qB}{c}\,\hat x\right)^2\right].
\label{eq:3.5.11}
\end{equation}
Since the operator $\hat p_y$ commutes with the Hamiltonian, we can
conclude that the momentum along the $y$ direction is a conserved
quantity and it can be replaced by its eigenvalue $\hbar k_y$. The
Hamiltonian is hence easily recognized as a shifted harmonic
oscillator Hamiltonian analogous to that discussed in the previous
example:
\[
\hat H=\frac{\hat p_x^2}{2m}
+\frac{1}{2}\,m\omega_c^2
\left(\hat x-x_B\,\hat 1\right)^2
+\frac{\hbar^2 k_y^2}{2m},
\]
\begin{equation}
\omega_c=\frac{|q|B}{mc},\qquad x_B=\frac{\hbar c k_y}{qB}.
\label{eq:3.5.12}
\end{equation}
The eigenfunctions of the above Hamiltonian can be written as
\begin{equation}
\Phi_{L,n}(x,y)=e^{ik_y y}\,\Phi_n(x-x_B),
\label{eq:3.5.13}
\end{equation}
which is essentially the product of a plane wave and a harmonic
oscillator eigenfunction, while the relevant eigenvalues read
\begin{equation}
E_n=\hbar\omega_c\!\left(n+\frac{1}{2}\right)
+\frac{\hbar^2 k_y^2}{2m}.
\label{eq:3.5.14}
\end{equation}
where the frequency $\omega_c$ is said the cyclotron frequency.

Before concluding this survey on applications of quantum oscillator
formalism, we discuss the case of Morse potential used to treat the
dynamical behavior of diatomic molecules and its approximation in
terms of harmonic oscillator potential. The Morse potential energy
function is of the form
\begin{equation}
V(r)=D_e\!\left(1-e^{-a(r-r_e)}\right)^2,
\label{eq:3.5.15}
\end{equation}
where $r$ is the distance between the atoms, $r_e$ the equilibrium
bond distance, $D_e$ is the well depth (defined relative to the
dissociated atoms) and $a$ accounts for the ``width'' of the
potential. The dissociation energy of the bond can be calculated by
subtracting the zero point energy $E(0)$ from the depth of the
well. By expanding $V(r)$ up to the lowest order in $(r-r_e)$, we
get
\begin{equation}
V(r)\approx\frac{a^2D_e}{2}\,(r-r_e)^2.
\label{eq:3.5.16}
\end{equation}
We can thus set $a^2D_e/2=\frac{1}{2}k$ and interpret this
approximation as a harmonic oscillator potential. The comparison
between the Morse potential (blue) and harmonic oscillator
potential (green) is shown in Fig.~\ref{fig:3.4}.

\begin{figure}[H]
\centering
\includegraphics[max width=0.95\linewidth,keepaspectratio]{fig3_4.png}\caption{Comparison between Morse (blue) and harmonic oscillator
(green) potentials.}
\label{fig:3.4}
\end{figure}

The Morse potential energy levels are not equally spaced as in the
case of the harmonic oscillator potential and their separation
decreases as the energy approaches the dissociation energy.

It is worth proving that the lowest order levels are well
approximated by
\begin{equation}
E_n\approx\hbar\omega\!\left(n+\frac{1}{2}\right),\qquad
\nu\equiv a\sqrt{\frac{2D_e}{m}},
\label{eq:3.5.17}
\end{equation}
while the corrections for higher orders read
\begin{equation}
E_n\approx\hbar\omega\!\left(n+\frac{1}{2}\right)
-\frac{\hbar^2\omega^2}{4D_e}\!\left(n+\frac{1}{2}\right)^2.
\label{eq:3.5.18}
\end{equation}

The use of the ordinary Perturbation Theory allows the derivation of
the corrections to the system wave functions.

% ------------------------------------------------------------------
\section{Coherent or Quasi-Classical States of Harmonic Oscillators}
\label{sec:3.6}

From eq.~\eqref{eq:3.4.10}, the position and momentum operators
can be written in terms of the creation--annihilation operators as
follows:
\begin{equation}
\hat\xi=\frac{1}{\sqrt{2}}\,(\hat a+\hat a^\dagger),\qquad
\hat p_\xi=\frac{i\hbar}{\sqrt{2}}\,(\hat a^\dagger-\hat a).
\label{eq:3.6.1}
\end{equation}

Let us now consider a Hamiltonian operator ruling a process of
dipole transitions between equally spaced quantum levels. Such an
operator can be cast in the form
\begin{equation}
\hat H=\hbar\Omega\,\hat\xi
\label{eq:3.6.2}
\end{equation}
and the associated Schr\"odinger equation can be written as
\begin{equation}
i\partial_t\Psi(\xi,t)=\frac{\Omega}{\sqrt{2}}\,(\hat a+\hat a^\dagger)\,\Psi(\xi,t),
\label{eq:3.6.3}
\end{equation}
which can be understood as a process in which a harmonic
oscillator level is created or destroyed as it happens during the
process of creation or destruction of photons in an electromagnetic
process. Eq.~\eqref{eq:3.6.3} admits the solution
\begin{equation}
\Psi(\xi,t)=e^{-i\Omega t(\hat a+\hat a^\dagger)/\sqrt{2}}\,\Psi(x,0),
\label{eq:3.6.4}
\end{equation}
i.e.\ assuming that the initial state is the vacuum, using the Dirac
notation, we can write
\begin{equation}
|\alpha(t)\rangle=
e^{-i\Omega t(\hat a+\hat a^\dagger)/\sqrt{2}}|0\rangle.
\label{eq:3.6.5}
\end{equation}
The explicit solution in closed form (see eq.~\eqref{eq:2.4.22}) can
be easily obtained
\noindent{\itshape (Hint: remind what we have learned about
operator ordering and disentanglement rules.)}
\begin{equation}
|\alpha(t)\rangle=
e^{-(\Omega t)^2/4}\,
e^{-i\Omega t\hat a^\dagger/\sqrt{2}}\,
e^{-i\Omega t\hat a/\sqrt{2}}\,|0\rangle.
\label{eq:3.6.6}
\end{equation}
which, on account of the fact that
\begin{equation}
e^{-i\Omega t\hat a/\sqrt{2}}\,|0\rangle=|0\rangle,
\label{eq:3.6.7}
\end{equation}
yields
\[
|\alpha(t)\rangle=
e^{-(\Omega t)^2/4}\,
\sum_{r=0}^{\infty}\frac{(-i\sqrt{2}\,\Omega t)^r}{2^{r/2}\,r!}\,
(\hat a^\dagger)^r|0\rangle
\]
\begin{equation}
=e^{-(\Omega t)^2/4}\,
\sum_{r=0}^{\infty}\frac{(-i\Omega t/\sqrt{2})^r}{r!}\,
|r\rangle.
\label{eq:3.6.8}
\end{equation}

The ``wavefunction'' $|\alpha(t)\rangle$ is given by a sum of
harmonic oscillator states and denotes a particular quantum state,
called \emph{coherent}, for the reasons we will specify below. The
probability that this state, at a given time $t$, be the $m$th
harmonic oscillator excited level, is
\begin{equation}
P_m(t)=|\langle m|\alpha(t)\rangle|^2
=e^{-(\Omega t)^2/2}\,\frac{(\Omega t/2)^{2m}}{m!}.
\label{eq:3.6.9}
\end{equation}
This expression, which is a consequence of the orthogonality
property~\eqref{eq:3.4.18}, is recognised to be a Poisson
distribution and the probabilities of finding quantum harmonic
oscillators excited levels $0,1,2,\ldots$ are shown in
Fig.~\ref{fig:3.5}.

\begin{figure}[H]
\centering
\includegraphics[max width=0.95\linewidth,keepaspectratio]{fig3_5.png}\caption{Probability of finding the coherent state in $|m\rangle$
for $m=0,1,2,\ldots$ as a function of $\Omega t$.}
\label{fig:3.5}
\end{figure}

We will just touch on the properties of coherent states which
appropriately describe the electromagnetic field of coherent sources
of laser type. One of the properties of these states is that they
are eigenstates of the annihilation operator. From the previous
equations, we find
\[
\hat a|\alpha(t)\rangle=
e^{-(\Omega t)^2/4}\,
\sum_{r=0}^{\infty}\frac{(-i\sqrt{2}\,\Omega t)^r}{2^{r/2}\,r!}\,
\hat a|r\rangle
=e^{-(\Omega t)^2/4}\,
\sum_{r=1}^{\infty}\frac{(-i\Omega t/\sqrt{2})^{r-1}}{(r-1)!}\,
|r-1\rangle
\]
\begin{equation}
=e^{-(\Omega t)^2/4}\,
\sum_{r=0}^{\infty}\frac{(-i\Omega t/\sqrt{2})^r}{r!}\,
|r\rangle
=c\,|\alpha(t)\rangle,\qquad
c=-\frac{i\Omega t}{\sqrt{2}}.
\label{eq:3.6.10}
\end{equation}
This apparently not significant result has striking consequences. It
is also worth stressing that the coherent state is not an
eigenvalue of the creation operator, i.e.
\begin{equation}
\hat a^\dagger|\alpha(t)\rangle\neq c^*|\alpha(t)\rangle.
\label{eq:3.6.11}
\end{equation}
On the other side, the identity
\begin{equation}
\langle\alpha(t)|\hat a^\dagger=\langle\alpha(t)|c^*
\label{eq:3.6.12}
\end{equation}
holds and, furthermore, it is also evident that
\begin{equation}
\langle\alpha(t)|\alpha(t)\rangle=1.
\label{eq:3.6.13}
\end{equation}

Let us now derive an important physical consequence from the
previous mathematical properties by proving that coherent states are
\emph{minimum uncertainty states}, i.e.\ they are quantum states
with minimum indetermination. To this aim, we evaluate the r.m.s.\
(root mean square) of position and momentum operators of coherent
states. According to the definition of r.m.s.\ and taking into
account the commutation rule~\eqref{eq:3.4.11}, one has
\begin{equation}
\sigma_\xi^2=
\langle\alpha(t)|\hat\xi^2|\alpha(t)\rangle-
\bigl(\langle\alpha(t)|\hat\xi|\alpha(t)\rangle\bigr)^2.
\label{eq:3.6.14}
\end{equation}
By using eq.~\eqref{eq:3.6.1} and other technicalities discussed
here, we find
\begin{equation}
\sigma_\xi^2=\frac{1}{2}\bigl(c^2+(c^*)^2+2|c|^2+1\bigr)
-\frac{1}{2}\bigl(c^2+(c^*)^2+2|c|^2\bigr)=\frac{1}{2},
\label{eq:3.6.15}
\end{equation}
and, similarly, for the momentum
\begin{equation}
\sigma_p^2=\frac{\hbar^2}{2}.
\label{eq:3.6.16}
\end{equation}
Going back to the dimensional counterparts, we get
\begin{equation}
\sigma_x^2=\frac{\hbar}{2m\omega},\qquad
\sigma_p^2=\frac{m\hbar\omega}{2},
\label{eq:3.6.17}
\end{equation}
and therefore
\begin{equation}
\sigma_x\,\sigma_p=\frac{\hbar}{2}.
\label{eq:3.6.18}
\end{equation}

On the other side the same procedure, applied to the $n$-states,
is shown that the uncertainty increases with the order of the level
\begin{equation}
(\sigma_x)_n^2=\frac{\hbar}{2m\omega}\,(2n+1),\qquad
(\sigma_p)_n^2=\frac{m\hbar\omega}{2}\,(2n+1),
\label{eq:3.6.19}
\end{equation}
\begin{equation}
(\sigma_x\,\sigma_p)_n=\frac{\hbar}{2}\,(2n+1).
\label{eq:3.6.20}
\end{equation}

This means that, with increasing $n$, the wave function is more
and more spread in both position and momentum. On the contrary, we
have found that coherent states have the peculiarity of being states
of minimum uncertainty. Among the quantum states, they are those
closer to a classical system.

This statement is however very qualitative and some further
comments are in order. In Fig.~\ref{fig:3.6}, we have reported a
comparison between classical and quantum probability for the
harmonic oscillators (see Rem.~1 for more details).

In the case of states with low $n$, there is a significant
difference between quantum and classical probabilities. It happens
that in the classical case, the oscillator will spend most of its
time far away from the position of equilibrium and this means that
the maximum probability is obtained for larger value of the
position. On the contrary, the quantum probability for $n=0$ is
maximum around $x=0$. We also note that there is a significant
portion of space which is not allowed classically. The quantum and
classical probabilities tend to become more similar with increasing
$n$, in full agreement with the correspondence principle.

It is now explained why a coherent state is closer to a quantum
state \noindent{\itshape (Hint: note that it is a superposition of
$n$-states\ldots)}.

\begin{exercise}
Reconsider the previous results, in particular their physical
meaning, and derive eqs.~\eqref{eq:3.6.19}.
\end{exercise}

\begin{figure}[H]
\centering
\includegraphics[max width=0.95\linewidth,keepaspectratio]{fig3_6.png}\caption{Comparison between classical and quantum probability for
the harmonic oscillators for (a)~$n=0$, (b)~$n=1$, (c)~$n=2$,
(d)~$n=3$, (e)~$n=10$.}
\label{fig:3.6}
\end{figure}

\medskip
\noindent{\itshape Remark~1.} A careful explanation is needed to
appreciate the physical content of Fig.~\ref{fig:3.6}. We have
mentioned the classical and quantum probability densities yielding
the probability of finding a particle at the position $x$ at the
time $t$. We have shown how to calculate the quantum probability
but the classical counterpart is less evident. We have reported a
very qualitative argument to evaluate the probability density by
noting that the particle spends more time in the region of space
where its velocity is lower. With reference to Fig.~\ref{fig:3.3}
where we have reported the quadratic potential, the probability
density $P(x)$ is maximum near the external points where the
velocity vanishes. The analytical derivation of $P(x)$ is given
below. We note that the probability of finding a particle of
crossing the position $\Delta x$ in the time interval $\Delta t$ is
\begin{equation}
P(x)\,\Delta x=2\,\frac{\Delta t}{T},
\label{eq:3.6.20}
\end{equation}
where $T$ is the oscillation period and the factor 2 is due to the
fact that the particle crosses the same region twice. It is
accordingly evident that
\begin{equation}
P(x)=\lim_{\Delta t\to0}\frac{2\,\Delta t}{T\,\Delta x}
=\frac{2}{T\,|\dot x|}.
\label{eq:3.6.21}
\end{equation}
It should now be reminded that for a classical harmonic
oscillator, we have the total energy of a particle with mass $m$
undergoing a quadratic potential to be
\begin{equation}
E=\frac{1}{2}\,m\dot x^2+\frac{1}{2}\,k\,x^2.
\label{eq:3.6.22}
\end{equation}
We can accordingly specify the quantities appearing in
eq.~\eqref{eq:3.6.21} as
\[
T=\frac{2\pi}{\omega},\qquad \omega=\sqrt{\frac{k}{m}},\qquad
|\dot x|=\omega\sqrt{a^2-x^2},\qquad
a^2=\frac{2E}{m\omega^2},
\]
thus getting
\begin{equation}
P(x)=\frac{1}{\pi\sqrt{a^2-x^2}}.
\label{eq:3.6.24}
\end{equation}
To make the comparison between classical and quantum
probabilities, we note that by replacing $E$ with the
corresponding quantum value, we get
\[
a^2=\sigma_x^2\,(2n+1),\qquad \sigma_x^2=\frac{\hbar}{2m\omega}.
\]
By using the quantity $a^2$, we obtain the reference dimensionless
position interval in which we compare the quantum and classical
probabilities.

% ------------------------------------------------------------------
\section{Jaynes--Cummings Model}
\label{sec:3.7}

In this section, we will introduce problems involving the combined
use of creation--annihilation operators and the Pauli matrices. We
will study the Jaynes--Cummings model (JC) which is the extension
of the already discussed two-level atom model (see Chapter~1) in
which a classical electromagnetic field ensures the coupling
between upper and lower levels.

If we do not make the hypothesis of a classical field, the analysis
of evolution of the system should be done by exploiting the
following Hamiltonian:
\begin{equation}
\hat H=\hat H_F+\hat H_I,\qquad
\hat H_F=\hbar\nu\,(\hat a^\dagger\hat a+\tfrac{1}{2}\hat\sigma_z),\qquad
\hat H_I=\hbar g\!\left(\hat a\,\hat\sigma_+\,e^{i\delta t}
+\hat a^\dagger\,\hat\sigma_-\,e^{-i\delta t}\right),
\label{eq:3.7.1}
\end{equation}
where $\nu$ is the frequency of the electromagnetic field and
$\omega$ is the separation between the two levels. The above
Hamiltonian should be understood as follows.
\begin{enumerate}
\item[(a)] The term $\hat H_F$ accounts for the free energy of the
      system, consisting of two contributions: the first associated
      with electromagnetic field (given by $\hbar\nu\hat a^\dagger\hat a$,
      related to the total number of photons $n_f$) and the second
      with the two-level part (given by $\frac{\hbar}{2}\omega\hat\sigma_z$).
      We emphasize that the operator $\hat\sigma_z$ is associated
      with $n_l$, i.e.\ the population difference between excited
      and ground states.
\end{enumerate}

\begin{figure}[H]
\centering
\includegraphics[max width=0.95\linewidth,keepaspectratio]{fig3_7.png}\caption{The Jaynes--Cummings model: coupling between a two-level
atom and a single radiation mode.}
\label{fig:3.7}
\end{figure}

\begin{exercise}
Prove that the total Hamiltonian commutes with the field part,
namely
\begin{equation}
[\hat H,\hat H_F]=0,
\label{eq:3.7.2}
\end{equation}
and that this corresponds to the following relation of conservation
\begin{equation}
n_f+n_l=k.
\label{eq:3.7.3}
\end{equation}
\end{exercise}

The previous ``law'' of conservation is extremely useful because it
allows the characterization of the dynamics of the system in terms
of a single quantum number which is the number of photons
exchanged during the interaction. We can therefore write the
physical states on which the JC Hamiltonian acts as follows
\begin{equation}
\Psi=\binom{|n+1\rangle}{|n\rangle},
\label{eq:3.7.4}
\end{equation}
(the upper level is indeed characterized by the presence of an extra
photon).

\begin{exercise}
Prove that the matrix element yielding the transition between two
states is just given by
\begin{equation}
\langle\Psi|\hat H|\Psi\rangle=
\begin{pmatrix}
\langle n|\hat H|n\rangle & \langle n|\hat H|n+1\rangle\\[4pt]
\langle n+1|\hat H|n\rangle & \langle n+1|\hat H|n+1\rangle
\end{pmatrix}
=
\frac{\hbar}{2}
\begin{pmatrix}
\omega & 2g\sqrt{n+1}\\
2g\sqrt{n+1} & -\omega
\end{pmatrix},
\label{eq:3.7.5}
\end{equation}
and that the energy eigenvalues characterizing upper and lower
levels are
\begin{equation}
E_\pm(n)=\hbar\!\left(n+\frac{1}{2}\right)\nu
\pm\frac{\hbar}{2}\,\Omega(n,\delta),\qquad
\Omega(n,\delta)=\sqrt{\delta^2+4g^2(n+1)},
\label{eq:3.7.6}
\end{equation}
where $\delta=\omega-\nu$. Find the link with the obtained results
in the case of a classical electromagnetic field.
\end{exercise}

% ------------------------------------------------------------------
\section{Classical Optics and Hermite Polynomials}
\label{sec:3.8}

In Chapter~2, we have discussed the solutions of the Helmholtz
equation, written in paraxial form, in terms of simple Gaussian
beams. Here we will prove that the solution of this equation is
compatible with a slightly more complicated form, involving Hermite
polynomials. These solutions will be called \emph{Hermite--Gauss}
beams.

Thus, we consider the heat equation
\begin{equation}
\partial_t F(x,t)=\alpha\,\partial_x^2 F(x,t),\qquad
F(x,0)=H_n(\beta x)\,e^{-x^2/(2\sigma^2)}.
\label{eq:3.8.1}
\end{equation}
Its solution can be obtained in terms of the Weierstrass transform
and reads
\begin{equation}
F_n(x,t)=\frac{1}{\sqrt{4\pi\alpha t}}
\int_{-\infty}^{\infty}
e^{-\frac{(x-\xi)^2}{4\alpha t}}\,
H_n(\beta\xi)\,e^{-\xi^2/(2\sigma^2)}\,d\xi.
\label{eq:3.8.2}
\end{equation}

The evaluation of the above integral is carried out by the use of
the generating function method, namely
\[
I(x,t;\lambda)=\sum_{n=0}^{\infty}\frac{\lambda^n}{n!}\,F_n(x,t)
=\frac{1}{\sqrt{4\pi\alpha t}}
 \int_{-\infty}^{\infty}
 e^{-\frac{(x-\xi)^2}{4\alpha t}}\,
 e^{2\lambda\beta\xi-\lambda^2}\,
 e^{-\xi^2/(2\sigma^2)}\,d\xi.
\]

The problem can therefore be easily solved using the rules of
integration of Gaussian functions. Therefore we find
\[
F_n(x,t)=
\frac{\sigma^n}{\sqrt{\sigma^2(t)}}\,
H_n\!\left(\frac{x}{\sigma(t)},\frac{4\alpha\beta^2 t}{\sigma^4(t)}\right)\,
\exp\!\left[-\frac{x^2}{2\sigma^2(t)}\right],
\]
\begin{equation}
\sigma(t)=\sqrt{\sigma^2+4\alpha t},
\label{eq:3.8.4}
\end{equation}
which, on account of eq.~\eqref{eq:3.3.9}, for
$\beta=\sigma^{-1}$, yields
\begin{equation}
F_n(x,t)=
\frac{1}{\sqrt{\sigma(t)}}\,
H_n\!\left(\frac{x}{\sigma(t)}\right)\,
\exp\!\left[-\frac{x^2}{2\sigma^2(t)}\right],
\label{eq:3.8.5}
\end{equation}
where $H_n$ is the canonical Hermite polynomial.

According to the previous result, we can write the solution of the
Helmholtz equation (2.5.2) as
\begin{equation}
F_{l,m}(r,z)=
F_0\!\left(\frac{W}{W(z)}\right)\,
H_l\!\left(\frac{\sqrt{2}\,x}{W(z)}\right)\,
H_m\!\left(\frac{\sqrt{2}\,y}{W(z)}\right)\,
e^{-\frac{x^2+y^2}{W^2(z)}}\,
e^{-i\varphi-i\frac{k}{2R(z)}(x^2+y^2)+i\zeta(z)(l+m+1)},
\label{eq:3.8.6}
\end{equation}
which represents the Gauss--Hermite optical modes (also called
$\mathrm{TEM}_{lm}$ modes) of the paraxial approximation. In
Fig.~\ref{fig:3.8}, we have reported the relevant structures of the
transverse distribution of the first few modes. It is also worth
noting that the emergence of these modes breaks the cylindrical
symmetry of the Gaussian beams, so they cannot be discussed within
the context of a scalar theory. This aspect of the problem cannot
however be discussed within the present book.

\begin{figure}[H]
\centering
\includegraphics[max width=0.95\linewidth,keepaspectratio]{fig3_8.png}\caption{Transverse distribution of the first Gauss--Hermite
$\mathrm{TEM}_{lm}$ modes.}
\label{fig:3.8}
\end{figure}

In this chapter, we have presented the theory and applications of
Hermite polynomials. We have used a non-conventional point of
view, with the merit of providing a straightforward and flexible
formalism. The reader is however invited to make a comparison
with the conventional (and more rigorous) treatment reported in
the bibliography below provided.

% ------------------------------------------------------------------
\section*{Bibliography}
\addcontentsline{toc}{section}{Bibliography}

\subsection*{Two Variable Hermite Polynomials}
\begin{enumerate}[label={[\arabic*]}]
\item P.~App\'el, J.~Kamp\'e de F\'eriet,
      \emph{Fonctions Hypcrg\'eom\'etriques et Hypersph\'eriques,
      polyn\^omes d'Hermite}, Gautier-Villars, Paris, 1926.
\item Digital Library of Mathematical Functions,
      \emph{Orthogonal Polynomials, Classical Orthogonal
      Polynomials, Sums}, (18.18), National Institute of
      Standards and Technology, retrieved 30 January 2015.
\item H.W.~Gould, A.T.~Hopper,
      ``Operational formulas connected with two generalizations
      of Hermite polynomials'',
      \emph{Duke Math.\ J.} \textbf{29}, pp.~51--63, 1962.
\item N.N.~Lebedev, \emph{Special Functions and Their
      Applications}, Dover Books on Mathematics, 1972.
\item G.~Dattoli, ``Generalized polynomials, operational
      identities and their applications'',
      \emph{J.\ Comput.\ Appl.\ Math.} \textbf{118},
      pp.~111--123, 2000.
\item G.~Dattoli, S.~Lorezutta, A.~Torre,
      ``Miscellaneous identities of generalized Hermite
      polynomials'', \emph{Matematiche (Catania)} \textbf{52},
      pp.~337--343, 1997.
\item G.~Dattoli, S.~Lorezutta, G.~Maino, A.~Torre,
      ``Theory of multi-index multivariable Bessel functions
      and Hermite polynomials'',
      \emph{Matematiche (Catania)} \textbf{52},
      pp.~179--197, 1997.
\item G.~Dattoli, A.~Torre, M.~Carpanese,
      ``Operational rules and arbitrary order Hermite
      generating functions'',
      \emph{J.\ Math.\ Anal.\ Appl.} \textbf{227},
      pp.~98--111, 1998.
\item M.~Artioli, G.~Dattoli,
      ``The Geometry of Hermite Polynomials'',
      Wolfram Demonstrations Project, Mar.~2015,
      \url{http://demonstrations.wolfram.com/
      TheGeometryOfHermitePolynomials}.
\end{enumerate}

\subsection*{Monomiality}
\begin{enumerate}[label={[\arabic*]},start=10]
\item G.~Dattoli, A.~Torre, G.~Mazzacurati,
      ``Quasi-monomials and isospectral problems'',
      \emph{Nuovo Cimento B} \textbf{112}, pp.~133--138, 1997.
\item G.~Dattoli, ``The Hermite and Laguerre Bessel Functions:
      A by Product of the Monomiality Principle'', in:
      D.~Cocolicchio, G.~Dattoli, H.M.~Srivastava (Eds.),
      \emph{Advanced Special Functions and Applications},
      Melfi, 1999, Aracne Editrice, Rome, p.~83, 2000.
\item P.~Blasiak, G.~Dattoli, A.~Horzela, K.~Penson,
      ``Representations of monomiality principle with
      Sheffer-type polynomials and boson normal ordering'',
      \emph{Phys.\ Lett.\ A} \textbf{352}, pp.~7--12, 2006.
\item J.F.~Steffensen, ``The poweroid, an extension of the
      mathematical notion of power'',
      \emph{Acta Math.} \textbf{73}, pp.~333--366, 1941.
\item A.V.~Turbiner, ``Quasi-exactly-solvable problems and
      $\mathfrak{sl}(2)$ algebra'',
      \emph{Commun.\ Math.\ Phys.} \textbf{118}, 46, 1988.
\end{enumerate}

\subsection*{Hermite Polynomials as Orthogonal Set}
\begin{enumerate}[label={[\arabic*]},start=15]
\item D.T.~Haimo, C.~Markett, ``A representation theory for
      solutions of a higher-order heat equation I'',
      \emph{J.\ Math.\ Anal.\ Appl.} \textbf{168},
      pp.~89--107, 1992.
\item D.T.~Haimo, C.~Markett, ``A representation theory for
      solutions of a higher-order heat equation II'',
      \emph{J.\ Math.\ Anal.\ Appl.} \textbf{168},
      pp.~289--305, 1992.
\item G.~Dattoli, B.~Germano, P.E.~Ricci,
      ``Comments on monomiality, ordinary polynomials and
      associated bi-orthogonal functions'',
      \emph{Appl.\ Math.\ Comput.} \textbf{154},
      pp.~219--227, 2004.
\item G.~Dattoli, H.M.~Srivastava, K.~Zhukovsky,
      ``Orthogonality properties of the Hermite and related
      polynomials'',
      \emph{J.\ Comput.\ Appl.\ Math.} \textbf{182}(1),
      pp.~165--172, 2005.
\item G.~Dattoli, B.~Germano, P.E.~Ricci,
      ``Higher order Hermite polynomials, associated
      bi-orthogonal functions and generalized heat
      equations'',
      \emph{Integral Transforms Spec.\ Funct.} \textbf{16},
      2005.
\end{enumerate}

\subsection*{Geometry of Hermite Polynomials}
\begin{enumerate}[label={[\arabic*]},start=20]
\item M.~Artioli, G.~Dattoli,
      ``The Geometry of Hermite Polynomials'',
      Wolfram Demonstrations Project, 2015,
      \url{http://demonstrations.wolfram.com/
      TheGeometryOfHermitePolynomials}.
\item M.~Artioli, G.~Dattoli,
      ``Geometric Properties of Generalized Hermite
      Polynomials'', Wolfram Demonstrations Project, 2015,
      \url{http://demonstrations.wolfram.com/
      GeometricPropertiesOfGeneralizedHermitePolynomials}.
\end{enumerate}

\subsection*{Applications}
\begin{enumerate}[label={[\arabic*]},start=22]
\item W.H.~Louisell, \emph{Quantum Statistical Properties of
      Radiation}, Wiley, New York, 1973.
\item G.~Dattoli, P.L.~Ottaviani, A.~Torre, L.~Vazquez,
      ``Evolution operator equations, integration with
      algebraic and finite difference methods: applications
      to physical problems in classical and quantum
      mechanics'', \emph{Riv.\ Nuovo Cimento} \textbf{20},
      pp.~1--133, 1997.
\end{enumerate}
